%! TeX root = ../main.tex

\subsection{Función de transferencia del sistema}
La función de transferencia $G(s)$ del sistema en el espacio de estados
se define de la forma

\begin{equation}
  G(s) = C(sI-A)^{-1}B + D
\end{equation}
 
\subsubsection{Desglose de $(sI-A)^{-1}$}

\begin{equation}
  \begin{aligned}
    sI-A &= \begin{bmatrix}
      s & 0 & 0 \\
      0 & s & 0 \\
      0 & 0 & s \\
    \end{bmatrix}
    - \begin{bmatrix}
      -\frac{R}{L}  & 0 &  -\frac{K_e}{L} \\
      0             & 0 &   1             \\
      \frac{K_t}{J} & 0 &  -\frac{B}{J}   \\
    \end{bmatrix}\\
    &= \begin{bmatrix}
      s+\frac{R}{L}    & 0-0 & 0+\frac{K_e}{L} \\
      0-0              & s-0 & 0-1             \\
      0-\frac{K_t}{J}  & 0-0 & s+\frac{B}{J}   \\
    \end{bmatrix}\\
    &= \begin{bmatrix}
      s+\frac{R}{L}  & 0 &  \frac{K_e}{L} \\
      0              & s & -1             \\
     -\frac{K_t}{J}  & 0 &  s+\frac{B}{J} \\
    \end{bmatrix}
  \end{aligned}
\end{equation}

Sabemos que la inversa de una matriz se define como en la ecuación \eqref{eq:InversaMatriz}

\begin{equation}
  Q^{-1} = \frac{1}{adj(Q)} Q^T
  \label{eq:InversaMatriz}
\end{equation}

Por lo tanto $(sI-A)^{-1} = \frac{1}{det(sI-A)}(sI-A)^T$ 

\begin{equation*}
  \begin{aligned}
    det(sI-A) &=
    \left(s+\frac{R}{L}\right)
    \begin{vmatrix}
      s & -1             \\
      0 &  s+\frac{B}{J} \\
    \end{vmatrix}\\
              &-
    0
    \begin{vmatrix}
      0             & -1             \\
      \frac{K_t}{J} &  s+\frac{B}{J} \\
    \end{vmatrix}
              +
    \frac{K_e}{L}
    \begin{vmatrix}
      0              & s  \\
     -\frac{K_t}{J}  & 0  \\
    \end{vmatrix}\\
    det(sI-A) &=
    \left(s+\frac{R}{L}\right)\left(s^2+\frac{B}{J}s\right) + 0 \\
    &+\frac{K_e}{L} \left(\frac{K_t}{J}s\right)\\
    det(sI-A) &=
    s^3 + \frac{B}{J}s^2 + \frac{R}{L}s^2 +\\ &+\frac{RB}{JL}s + \frac{K_t K_e}{JL}s\\
  \end{aligned}
\end{equation*}

\begin{equation}
  \begin{aligned}
    &det(sI-A) =\\
    &s^3 + \left(\frac{B}{J} + \frac{R}{L}\right)s^2 
    + \left(\frac{RB + K_t K_e}{JL}\right)s\\
  \end{aligned}
\end{equation}

\begin{equation*}
  \begin{aligned}
    &cof(sI-A) =\\ &\begin{bmatrix}
%%%%%%%%%TOP
      (-1)^{1+1}
      \begin{vmatrix}
        s             & -1              \\
        0             &  s+\frac{B}{J}  \\
      \end{vmatrix}&
      (-1)^{1+2}
      \begin{vmatrix}
        0             & -1              \\
       -\frac{K_t}{J} &  s+\frac{B}{J}  \\
      \end{vmatrix}&
      (-1)^{1+3}
      \begin{vmatrix}
        0             & s               \\
       -\frac{K_t}{J} & 0               \\
      \end{vmatrix}\\
%%%%%%%%MID
      (-1)^{2+1}
      \begin{vmatrix}
        0             &  \frac{K_e}{L} \\
        0             &  s+\frac{B}{J} \\
      \end{vmatrix}&
      (-1)^{2+2}
      \begin{vmatrix}
        s+\frac{R}{L} &  \frac{K_e}{L} \\
       -\frac{K_t}{J} &  s+\frac{B}{J} \\
      \end{vmatrix}&
      (-1)^{2+3}
      \begin{vmatrix}
        s+\frac{R}{L} & 0              \\
       -\frac{K_t}{J} & 0              \\
      \end{vmatrix}\\
%%%%%%%%BTM
      (-1)^{3+1}
      \begin{vmatrix}
        0 &  \frac{K_e}{L}             \\
        s & -1                         \\
      \end{vmatrix}&
      (-1)^{3+2}
      \begin{vmatrix}
        s+\frac{R}{L} &  \frac{K_e}{L} \\
        0             & -1             \\
      \end{vmatrix}&
      (-1)^{3+3}
      \begin{vmatrix}
        s+\frac{R}{L} & 0              \\
        0             & s              \\
      \end{vmatrix}\\
    \end{bmatrix}\\
  \end{aligned}
\end{equation*}

\begin{equation}
  \begin{aligned}
    &cof(sI-A) =\\
    &\begin{bmatrix}
      s^2 + \frac{B}{J}s  &  \frac{K_t}{J}                                                  &  \frac{K_t}Js   \\
      0                   &  s^2 + \left(\frac{R}L + \frac{B}J\right)s + \frac{K_e K_t}{JL} &  0              \\
      -\frac{K_e}L s      &  s +\frac{R}L                                                   &  s^2+\frac{R}Ls \\
    \end{bmatrix}
  \end{aligned}
\end{equation}

\begin{equation}
  \begin{aligned}
    &adj(sI-A) =\\
    &\begin{bmatrix}
      s^2 + \frac{B}{J}s                                             &   
      0                                                              &   
     -\frac{K_e}L s                                                  \\  
      \frac{K_t}{J}                                                  & 
      s^2 + \left(\frac{R}L + \frac{B}J\right)s + \frac{K_e K_t}{JL} & 
      s +\frac{R}L                                                   \\
      \frac{K_t}Js                                                   &
      0                                                              &
      s^2+\frac{R}Ls                                                   \\
    \end{bmatrix}
  \end{aligned}
\end{equation}

\subsubsection{Función de transferencia del sistema}

\begin{equation*}
  q(s) = s^2 + \left(\frac{R}L + \frac{B}J\right)s + \frac{K_e K_t}{JL}
\end{equation*}

\begin{multline*}
  G(s) =
    \frac{1}{s^3 + \left(\frac{B}{J} + \frac{R}{L}\right)s^2 + \left(\frac{RB + K_t K_e}{JL}\right)s}\\
    \begin{bmatrix}
      0 & 0 & 1 \\
    \end{bmatrix}
    \begin{bmatrix}
      s^2 + \frac{B}{J}s  &   
      0                   &   
     -\frac{K_e}L s       \\  
      \frac{K_t}{J}       & 
      q(s)                & 
      s +\frac{R}L        \\
      \frac{K_t}Js        &
      0                   &
      s^2+\frac{R}Ls      \\
    \end{bmatrix}
    \begin{bmatrix}
      \frac1L \\ 0 \\ 0 \\
    \end{bmatrix}
\end{multline*}

\begin{multline*}
  G(s) =
    \frac{1}{s^3 + \left(\frac{B}{J} + \frac{R}{L}\right)s^2 + \left(\frac{RB + K_t K_e}{JL}\right)s}\\
    \frac1L
    \begin{bmatrix}
      0 & 0 & 1 \\
    \end{bmatrix}
    \begin{bmatrix}
      s^2 + \frac{B}{J}s  \\
      \frac{K_t}{J}       \\
      \frac{K_t}Js        \\
    \end{bmatrix}
\end{multline*}

\begin{equation}
  G(s) =
    \frac{\frac{K_t}{JL}s}{s^3 + \left(\frac{B}{J} + \frac{R}{L}\right)s^2 + \left(\frac{RB + K_t K_e}{JL}\right)s}
\end{equation}


\subsubsection{Función de transferencia con valores experimentales}

A partir de los mismos valores experimentales: $R = 2.1\ \Omega$, $L = 5\times10^{-3}\ \text{H}$, $K_t = 0.12\ \text{N·m/A}$, $K_e = 0.12\ \text{V/(rad/s)}$, $J = 1\times10^{-4}\ \text{kg·m}^2$ y $B = 1\times10^{-5}\ \text{N·m·s/rad}$, la función de transferencia numérica resulta:

\begin{equation*}
  \begin{aligned}
    &G(s) =\\
    &\frac{\tfrac{0.12}{1\mathrm{e}{-4} \cdot 5\mathrm{e}{-3}}s}{s^{3} + \left(\tfrac{1\mathrm{e}{-5}}{1\mathrm{e}{-4}} + \tfrac{2.1}{5\mathrm{e}{-3}}\right)s^{2} + \left(\tfrac{2.1\cdot1\mathrm{e}{-5} + 0.12\cdot0.12}{1\mathrm{e}{-4} \cdot 5\mathrm{e}{-3}}\right)s}\\
    &G(s) =
    \frac{2.4\mathrm{e}^{5}\,s}{s^{3} + 4.201\mathrm{e}^{2}\,s^{2} + 2.8842\mathrm{e}^{4}\,s} 
  \end{aligned}
\end{equation*}

\begin{equation}
G(s) = \frac{2.4\mathrm{e}^{5}}{s^{2} + 4.201\mathrm{e}^{2}\,s + 2.8842\mathrm{e}^{4}}
\label{eq:TransferFunctionNumeric}
\end{equation}
